\chapter{多体系の量子 Wasserstein 距離 \texorpdfstring{$W_1$}{W1}}
\label{ch:quantum-w1}

本章では De Palma--Marvian--Trevisan--Lloyd
\emph{The Quantum Wasserstein Distance of Order 1}
（arXiv:2009.04469）の量子 $W_1$ を導入する．この定義は $n$-qubit 系の tensor 積構造，
すなわち「どの場所を変更したか」を幾何として使う．

\section{局所的な差}

以下

\[
  \mathcal{H}_n=(\C^2)^{\otimes n},
  \qquad
  \mathcal{O}_n^0
  :=\{X\in\Obs{\mathcal{H}_n}\mid\Tr X=0\}
\]

とおく．二つの量子状態の差は必ず $\mathcal{O}_n^0$ に属する．

\begin{definition}{隣接する量子状態}{q-neighboring}
  $\rho,\sigma\in\States{\mathcal{H}_n}$ が\textbf{隣接する}とは，ある $i\in\{1,\ldots,n\}$ で

  \[
    \Tr_i\rho=\Tr_i\sigma
  \]

  となることをいう．すなわち，第 $i$ qubit を捨てれば二状態を区別できない．
\end{definition}

一般の $X\in\mathcal{O}_n^0$ を，一つの場所だけに局所化された差の和へ分解する．

\begin{definition}{量子 $W_1$ ノルム}{q-w1-norm}
  $X\in\mathcal{O}_n^0$ に対し

  \[
    \norm{X}_{W_1}
    :=\frac12
    \min\left\{
      \sum_{i=1}^n\norm{X^{(i)}}_1
      \;\middle|\;
      X=\sum_{i=1}^nX^{(i)},\quad
      X^{(i)}\in\mathcal{O}_n^0,\quad
      \Tr_iX^{(i)}=0
    \right\}
  \]

  と定める．
\end{definition}

まず，この最小化が常に意味をもつことを確認する．

\begin{lemma}{局所分解の存在}{q-local-decomposition}
  任意の $X\in\mathcal{O}_n^0$ に対し，定義~\ref{def:q-w1-norm}の制約をみたす
  $(X^{(i)})_{i=1}^n$ が存在し，最小値は達成される．
\end{lemma}

\begin{proof}
  第 $i$ qubit を最大混合状態で置き換える線形写像を

  \[
    \mathcal{E}_i(X):=\frac{\Id_2}{2}\otimes\Tr_iX
  \]

  と書く．tensor 因子は元の順序に戻すものとする．異なる $\mathcal{E}_i$ は可換であり，
  $\Tr_i\mathcal{E}_i(X)=\Tr_iX$ である．
  $\mathcal{E}_{1:i}:=\mathcal{E}_1\cdots\mathcal{E}_i$，
  $\mathcal{E}_{1:0}:=\Id$ とおき

  \[
    X^{(i)}:=\mathcal{E}_{1:i-1}(X)-\mathcal{E}_{1:i}(X)
  \]

  と定めると，$\Tr_iX^{(i)}=0$ である．また
  $\mathcal{E}_{1:n}(X)=(\Tr X)\Id_{2^n}/2^n=0$ だから，望む telescope 和

  \[
    \sum_{i=1}^nX^{(i)}=X-\mathcal{E}_{1:n}(X)=X
  \]

  を得る．よって許容集合は空でない．

  最小化列は目的関数がある定数以下となるように取れる．すると各 $\norm{X^{(i)}}_1$ が
  一様に有界である．有限次元なので収束部分列が取れ，線形な制約は極限でも保たれる．
  跡ノルムの連続性から極限が最小値を達成する．
\end{proof}

\begin{proposition}{隣接状態による表示}{q-neighbor-form}
  定義~\ref{def:q-w1-norm}は

  \[
    \norm{X}_{W_1}
    =\min\left\{
      \sum_{i=1}^nc_i
      \;\middle|\;
      X=\sum_{i=1}^nc_i(\rho^{(i)}-\sigma^{(i)}),\quad
      c_i\geq0,\quad
      \Tr_i\rho^{(i)}=\Tr_i\sigma^{(i)}
    \right\}
  \]

  と同値である．ここで $\rho^{(i)},\sigma^{(i)}\in\States{\mathcal{H}_n}$ である．
\end{proposition}

\begin{proof}
  許容な $X^{(i)}$ の Jordan 分解を
  $X^{(i)}=X^{(i)}_+-X^{(i)}_-$ とし，
  $c_i:=\norm{X^{(i)}}_1/2$ とおく．$c_i>0$ なら

  \[
    \rho^{(i)}:=X^{(i)}_+/c_i,
    \qquad
    \sigma^{(i)}:=X^{(i)}_-/c_i
  \]

  は状態である．さらに $\Tr_iX^{(i)}=0$ より両者の部分跡は等しい．
  逆に $X^{(i)}=c_i(\rho^{(i)}-\sigma^{(i)})$ とおけば
  $\Tr_iX^{(i)}=0$ かつ $\norm{X^{(i)}}_1\leq2c_i$ である．
  両方向の表示に対して最小値を取れば等号を得る．
\end{proof}

\section{ノルムと距離性}

\begin{lemma}{跡ノルムによる下界}{q-w1-lower-bound}
  任意の $X\in\mathcal{O}_n^0$ に対し

  \[
    \frac12\norm{X}_1\leq\norm{X}_{W_1}
  \]

  が成り立つ．
\end{lemma}

\begin{proof}
  任意の許容分解について跡ノルムの三角不等式から

  \[
    \norm{X}_1
    =\norm{\sum_iX^{(i)}}_1
    \leq\sum_i\norm{X^{(i)}}_1
  \]

  である．右辺の半分について最小値を取ればよい．
\end{proof}

\begin{theorem}{量子 $W_1$ はノルムである}{q-w1-is-norm}
  $X\mapsto\norm{X}_{W_1}$ は実線形空間 $\mathcal{O}_n^0$ 上のノルムである．
\end{theorem}

\begin{proof}
  有限性は補題~\ref{lem:q-local-decomposition}による．定義から非負であり，
  補題~\ref{lem:q-w1-lower-bound}より
  $\norm{X}_{W_1}=0$ なら $\norm{X}_1=0$，したがって $X=0$ である．

  $a\in\R$ と許容分解 $(X^{(i)})_i$ に対し $(aX^{(i)})_i$ は $aX$ の許容分解だから
  $\norm{aX}_{W_1}\leq\abs{a}\norm{X}_{W_1}$ である．$a\neq0$ のとき $1/a$ を
  $aX$ に適用すれば逆向きの不等式を得る．$a=0$ は明らかである．

  最後に $X,Y$ の最小分解をそれぞれ $(X^{(i)})_i,(Y^{(i)})_i$ とすると，
  $(X^{(i)}+Y^{(i)})_i$ は $X+Y$ の許容分解である．よって

  \[
    \norm{X+Y}_{W_1}
    \leq\frac12\sum_i\norm{X^{(i)}+Y^{(i)}}_1
    \leq\norm{X}_{W_1}+\norm{Y}_{W_1}.
  \]

  以上で正定値性，絶対斉次性，三角不等式が示された．
\end{proof}

\begin{definition}{量子 Wasserstein 距離 $W_1$}{q-w1-distance}
  $\rho,\sigma\in\States{\mathcal{H}_n}$ に対し

  \[
    \Wass_1^{\rm q}(\rho,\sigma):=\norm{\rho-\sigma}_{W_1}
  \]

  と定める．
\end{definition}

\begin{corollary}{量子 $W_1$ の距離性}{q-w1-metric}
  $\Wass_1^{\rm q}$ は $\States{\mathcal{H}_n}$ 上の距離である．
\end{corollary}

\begin{proof}
  非負性と対称性はノルムから従う．また
  $\Wass_1^{\rm q}(\rho,\sigma)=0$ なら
  $\rho-\sigma=0$ だから $\rho=\sigma$ であり，逆は明らかである．三角不等式は

  \[
    \norm{\rho-\tau}_{W_1}
    \leq\norm{\rho-\sigma}_{W_1}+\norm{\sigma-\tau}_{W_1}
  \]

  そのものである．
\end{proof}

\begin{example}{一つの qubit}{q-one-qubit}
  $n=1$ では許容分解は $X^{(1)}=X$ しかないので

  \[
    \Wass_1^{\rm q}(\rho,\sigma)=\frac12\norm{\rho-\sigma}_1
    =D_{\rm tr}(\rho,\sigma).
  \]

  多体系では，各場所に局所化した差のコストを足すため，単なる跡距離とは異なる幾何になる．
\end{example}
